% Submit only if NSF/Georgia Tech classify a budget-supported project role as a
% CI Professional. The current planning budget has no separate RSE line.
\section*{CI Professional Mentoring and/or Professional Development Plan}

\subsection*{Covered Role and Individual Development Plan}

If the PI's budget-supported software-engineering role is classified as a CI
Professional, this plan applies to that effort; if a research software
engineer or programmer is added, the same activities will apply and the final
document will name the covered role(s). Within the first month, each covered
individual will prepare a development plan with the project science lead and
an independent Georgia Tech CI/research-software mentor. The plan will identify
technical, scientific, communication, leadership, and career objectives and
will be reviewed semiannually.

\subsection*{Mentoring and Professional Development}

The project will provide protected time and project-supported opportunities
for the following activities:

\begin{itemize}[leftmargin=*]
  \item \textbf{Research-software leadership:} ownership of architectural
  decision records, release planning, code review, incident retrospectives,
  maintenance estimates, and contributor onboarding, with quarterly feedback
  from the PI/Co-PI or independent mentor as appropriate.
  \item \textbf{Secure and reliable CI:} training in threat modeling,
  least-privilege remote execution, dependency/SBOM practices, vulnerability
  response, backups, and incident response. The CI Professional will
  participate in an external security or science-gateway review and lead the
  closure of at least one resulting recommendation.
  \item \textbf{Scientific collaboration:} monthly sessions with the science
  lead on model assumptions, numerical validation, uncertainty, and
  reproducibility; pairing with a PISM user/developer on adapter acceptance;
  and practice communicating engineering limitations to domain scientists.
  \item \textbf{Scholarship and visibility:} co-authorship or credited
  contributions on software/data papers and release records; at least one
  technical presentation, poster, tutorial, or workshop leadership role per
  year; and support to prepare abstracts, manuscripts, and future proposals.
  \item \textbf{Career development:} annual review of career pathways,
  promotion criteria, portfolio/CV, leadership goals, professional networks,
  and relevant training. Up to one focused course or conference per year may
  be selected within available project/institutional resources.
  \item \textbf{Responsible professional practice:} training in research
  integrity, accessibility, inclusive collaboration, software/data licensing,
  privacy-preserving metrics, credit and authorship, and responsible disclosure
  of security vulnerabilities.
\end{itemize}

\subsection*{Evidence and Accountability}

Semiannual reviews will record completed training, new responsibilities,
products, presentations, mentoring received/provided, and next objectives.
The project will assess whether the individual gains independent ownership of
a production component, demonstrates cross-disciplinary leadership, and has a
documented advancement path. Activities and outcomes will be summarized in
annual reports without disclosing private career discussions. Mentoring is a
project responsibility and will not depend solely on informal availability.
